\documentclass[11pt]{article}

\usepackage[margin=1.1in]{geometry}
\usepackage{amsmath,amssymb,amsthm,mathtools}
\usepackage{enumitem}
\usepackage{hyperref}
\usepackage{microtype}
\usepackage{booktabs}
\usepackage{mathrsfs}

\hypersetup{
  colorlinks=true,
  linkcolor=blue,
  citecolor=blue,
  urlcolor=blue,
  hypertexnames=false,
  pdftitle={Microscopic Cycles and Sparse Entries for Regular Luce Profiles},
  pdfauthor={Christopher D. Long}
}

\newtheorem{theorem}{Theorem}[section]
\newtheorem{proposition}[theorem]{Proposition}
\newtheorem{lemma}[theorem]{Lemma}
\newtheorem{corollary}[theorem]{Corollary}
\newtheorem{conjecture}[theorem]{Conjecture}
\theoremstyle{remark}
\newtheorem{remark}[theorem]{Remark}
\newtheorem{warning}[theorem]{Warning}
\newtheorem{definition}[theorem]{Definition}

\newcommand{\Pp}{\mathbb P}
\newcommand{\Ee}{\mathbb E}
\newcommand{\Rr}{\mathbb R}
\newcommand{\Nn}{\mathbb N}
\newcommand{\Zz}{\mathbb Z}
\newcommand{\one}{\mathbf 1}
\newcommand{\TV}{\mathrm{TV}}
\newcommand{\Pois}{\operatorname{Pois}}
\newcommand{\Exp}{\operatorname{Exp}}
\newcommand{\PPP}{\operatorname{PPP}}
\newcommand{\Var}{\operatorname{Var}}
\newcommand{\Cov}{\operatorname{Cov}}
\newcommand{\diag}{\operatorname{diag}}
\newcommand{\dd}{\,d}
\newcommand{\eps}{\varepsilon}
\newcommand{\calD}{\mathcal D}
\newcommand{\calC}{\mathcal C}
\newcommand{\calE}{\mathcal E}
\newcommand{\calG}{\mathcal G}

\title{Microscopic Cycles and Sparse Entries for Regular Luce Profiles}
\author{Christopher D. Long}
\date{Audited current manuscript --- localized clarifications only}

\begin{document}
\maketitle

\begin{abstract}
Let a Luce permutation be generated by independent exponential clocks with
rates
\[
 \theta_{n,i}=a_n\{f(i/n)+\varepsilon_{n,i}\},
 \qquad \max_i|\varepsilon_{n,i}|\longrightarrow0,
\]
where \(f\) is continuous and bounded above and below by positive constants.
This manuscript proves joint independent Poisson limits for every fixed vector
of cycle counts, identifies their means as traces of the microscopic Luce
kernel, and proves exponential convergence of those means to the uniform
values \(1/r\).  The same local-entry theory yields fixed-width diagonal-band
laws, sparse-entry Poisson limits, thin-graph laws, and overlap limits with
deterministic permutations.  The proof includes the separated bulk estimate
and the endpoint, collision, and target-crowding envelopes needed to sum fixed
factorial moments globally.
\end{abstract}

\tableofcontents

\section{Status and provenance}

This standalone source is a mechanical extraction from the frozen
139-page manuscript at repository commit
\texttt{6178a99a8f2ee7b6416edd026e7780b9b8fa9fa0}.  The mathematical text
between the archival sentinels below is copied byte-for-byte from the frozen
source.  Only the standalone preamble, abstract, this provenance notice,
section wrapper, and selected bibliography have been added.

The initial extraction preserved the profile hypothesis, all equality and
limit statements, every endpoint and collision envelope used by the
factorial-moment arguments, and all internal labels and references.  A later
adversarial audit found no fatal gap and recommended six localized
clarifications.  The present source differs from the frozen mathematical
blocks only by the marked clarification passages below; no hypothesis,
conclusion, limiting constant, or mode of convergence has been strengthened.
The audit and exact repair record are kept under \texttt{audits/}.

\section{Preliminaries}

% BEGIN ARCHIVED BLOCK: exponential-race-and-gumbel-input
\subsection{The exponential race}
Let \(\theta_1,\ldots,\theta_n>0\), and let \(X_i\sim\Exp(\theta_i)\) be independent.  For distinct labels \(i_1,\ldots,i_n\),
\[
\Pp(X_{i_1}<\cdots<X_{i_n})
= \prod_{r=1}^n
 \frac{\theta_{i_r}}{\theta_{i_r}+\theta_{i_{r+1}}+\cdots+\theta_{i_n}}.
\]
This is the standard exponential race identity and follows from the memoryless property.


\subsection{Gumbel priorities}\label{subsec:gumbel-priorities}
A standard Gumbel random variable has distribution
\[
        \Pp(G\le y)=\exp(-e^{-y}),
        \qquad y\in\Rr .
\]
If \(E_i\sim\Exp(1)\) are independent, then \(-\log E_i\) are independent standard Gumbels.  Since
\[
        X_i\stackrel{d}=E_i/\theta_i,
\]
the exponential race may equivalently be generated by the shifted priorities
\[
        Y_i=-\log X_i=\log\theta_i+G_i,
        \qquad 1\le i\le n,
\]
ordered in decreasing order.  Thus
\[
        \sigma(i)=1+\#\{j:Y_j>Y_i\}.
\]
For later reference we record the exact threshold formulas.  The density and upper-tail function of \(Y_i\) are
\[
        f_i(y)=\theta_i e^{-y}\exp(-\theta_i e^{-y}),
        \qquad
        \Pp(Y_i>y)=1-\exp(-\theta_i e^{-y}).
\]

\begin{lemma}[Gumbel threshold identities]\label{lem:gumbel-threshold-identities}
For a single specified label \(p\),
\[
        \Pp(\sigma(p)=q)
        =\int_{-\infty}^{\infty}
        f_p(y)\,
        \Pp\left(\sum_{j\ne p}\one_{\{Y_j>y\}}=q-1\right)\dd y .
\]
More generally, let \(p_1,\ldots,p_\ell\) be distinct labels and let
\(q_1<\cdots<q_\ell\) be distinct target positions.  Put
\(P=\{p_1,\ldots,p_\ell\}\) and
\[
        S_a(y)=\sum_{j\notin P}\one_{\{Y_j>y_a\}},
        \qquad 1\le a\le\ell .
\]
Then
\[
\begin{aligned}
&\Pp(\sigma(p_a)=q_a,\ 1\le a\le\ell) \\
&\quad=\int_{y_1>\cdots>y_\ell}
        \prod_{a=1}^\ell f_{p_a}(y_a)\,
        \Pp(S_a(y)=q_a-a,\ 1\le a\le\ell)
        \dd y_1\cdots\dd y_\ell .
\end{aligned}
\]
\end{lemma}

\begin{proof}
Condition on the specified priorities.  For one label, the event
\(\sigma(p)=q\) is exactly the event that precisely \(q-1\) nonspecified
priorities exceed the observed value of \(Y_p\).  This gives the first
identity.

For the multi-point identity, the ordering of the target positions forces the
specified priorities to satisfy \(Y_{p_1}>\cdots>Y_{p_\ell}\).  Conditional on
\(Y_{p_a}=y_a\) in this chamber, exactly \(a-1\) specified priorities are above
\(y_a\).  Hence \(\sigma(p_a)=q_a\) is equivalent to
\(S_a(y)=q_a-a\), simultaneously for \(1\le a\le\ell\).  Integrating against the
independent specified-priority densities gives the formula.
\end{proof}

\begin{remark}[The derivative cancellation]\label{rem:gumbel-derivative-cancellation}
The one-point identity explains the local-entry density used later.  Let
\[
        M_{n,P}(y)=\sum_{j\notin P}\Pp(Y_j>y)
        =\sum_{j\notin P}(1-e^{-\theta_j e^{-y}}).
\]
At a rank saddle \(y_*\), where \(M_{n,P}(y_*)\) equals the target rank up to
\(O(1)\), a one-dimensional local CLT gives formally
\[
        \Pp(\sigma(p)=q)
        \approx
        \frac{f_p(y_*)}{|M_{n,P}'(y_*)|}.
\]
Here
\[
        |M_{n,P}'(y_*)|
        =\sum_{j\notin P}\theta_j e^{-y_*}e^{-\theta_j e^{-y_*}},
\]
while
\[
        f_p(y_*)=\theta_p e^{-y_*}e^{-\theta_p e^{-y_*}}.
\]
Thus the common factor \(e^{-y_*}\) cancels.  In the regular-profile notation of
Section \ref{sec:regular}, the exact finite-$n$ saddle uses
$A_{n,P}(t)=(q-1)/n$.  Replacing $(q-1)/n$ by $q/n$ changes only a lower-order
term, and putting \(t=e^{-y_*}\) gives
\[
        \frac{f_p(y_*)}{|M_{n,P}'(y_*)|}
        =\frac1n\frac{\theta_p e^{-\theta_p t}}{B_{n,P}(t)}
        \longrightarrow
        \frac1n\rho(p/n,q/n).
\]
This is the Gumbel form of the exponential-clock saddle calculation in
Proposition \ref{prop:bulk-local-final}.
\end{remark}

% END ARCHIVED BLOCK: exponential-race-and-gumbel-input

% BEGIN ARCHIVED BLOCK: analytic-input
\subsection{Analytic input}
We use standard estimates for sums of independent Bernoulli variables.  They are stated explicitly to make clear what is invoked.

\begin{lemma}[Point probability and Bernstein bounds]\label{lem:bernstein-point}
Let \(Y=\sum_{j=1}^N \xi_j\), where the \(\xi_j\) are independent, not necessarily identically distributed, Bernoulli variables, and set \(\mu=\Ee Y\), \(v=\Var Y\).  There is an absolute constant \(C\) such that
\[
        \sup_m \Pp(Y=m)\le \frac{C}{\sqrt{1+v}}.
\]
Moreover, there is an absolute constant \(c>0\) such that, for every \(u\ge0\),
\[
        \Pp(|Y-\mu|\ge u)\le
        2\exp\left[-c\min\left(\frac{u^2}{v+1},u\right)\right].
\]
\end{lemma}

\begin{proof}
The point bound is the classical concentration-function estimate for sums of Bernoulli variables.  It follows from Fourier inversion and the standard inequality
\[
        |1-p+pe^{it}|
        \le \exp\left[-\frac{2}{\pi^2}p(1-p)\operatorname{dist}(t,2\pi\mathbb Z)^2\right],
\]
where \(\operatorname{dist}(t,2\pi\mathbb Z)\) denotes distance from \(t\) to the nearest multiple of \(2\pi\).  Bernstein's inequality gives the second bound for bounded independent summands.  See Petrov \cite[Chapter VII]{Petrov1975} for both estimates.
\end{proof}

\begin{lemma}[Uniform lattice local CLT with a tilted local bound]
\label{lem:llt}
Fix a dimension $d$ and a support bound $B<\infty$.  Let
\[
        Y_{n,m}=\sum_{j=1}^{N_n}\zeta_{n,m,j},
        \qquad m\in\mathcal M_n,
\]
be sums of independent $\mathbb Z^d$-valued random vectors satisfying
$\|\zeta_{n,m,j}\|\le B$.  Write
\[
        \mu_{n,m}=\Ee Y_{n,m},
        \qquad \Sigma_{n,m}=\Var(Y_{n,m}),
\]
and suppose that, uniformly in $m$,
\[
        c s_n I_d\le\Sigma_{n,m}\le C s_n I_d,
        \qquad s_n\longrightarrow\infty.
\]
For $\lambda\in\mathbb R^d$, let $\Pp_\lambda$ denote the product exponential
tilt
\[
 \Pp_\lambda(\zeta_{n,m,j}=x)
 =\frac{e^{\lambda\cdot x}\Pp(\zeta_{n,m,j}=x)}
 {\Ee e^{\lambda\cdot\zeta_{n,m,j}}}.
\]
Assume that there is $\lambda_0>0$ for which the following two properties hold
uniformly for $\|\lambda\|\le\lambda_0$:
\begin{enumerate}[label=(\roman*)]
\item the tilted covariance matrices satisfy
\[
        c_0s_nI_d\le\Var_\lambda(Y_{n,m})\le C_0s_nI_d;
\]
\item for every $\varepsilon>0$ there are $\kappa_\varepsilon>0$ and
$n_\varepsilon$ such that
\[
 \left|\prod_{j=1}^{N_n}
 \Ee_\lambda e^{iu\cdot\zeta_{n,m,j}}\right|
 \le e^{-\kappa_\varepsilon s_n}
\]
whenever $n\ge n_\varepsilon$ and
$u\in[-\pi,\pi]^d$ has distance at least $\varepsilon$ from
$2\pi\mathbb Z^d$.
\end{enumerate}
Then, uniformly in $m$ and $z\in\mathbb Z^d$,
\begin{align*}
 \Pp(Y_{n,m}=z)
 &=(2\pi)^{-d/2}(\det\Sigma_{n,m})^{-1/2}\\
 &\quad\times
 \exp\left[-\frac12(z-\mu_{n,m})^T
 \Sigma_{n,m}^{-1}(z-\mu_{n,m})\right]
 +o(s_n^{-d/2}).
\end{align*}
Moreover,
\[
        \sup_z\Pp(Y_{n,m}=z)\le C_1s_n^{-d/2}.
\]
For every fixed $K<\infty$ and $\gamma<1/2$, there are $c_1,C_2>0$ such
that, uniformly when
\[
        \|z-\mu_{n,m}\|\le K s_n^{1/2+\gamma},
\]
one has
\[
 \Pp(Y_{n,m}=z)
 \le C_2s_n^{-d/2}
 \exp\left[-c_1\frac{\|z-\mu_{n,m}\|^2}{s_n}\right].
\]
All constants depend only on the fixed dimension, the support bound, and the
uniform constants in the hypotheses.
\end{lemma}

\begin{proof}
We suppress $n,m$ and write $Y,\mu,\Sigma,s$.  Let
$X_j=\zeta_j-\Ee\zeta_j$ and
\[
        \Psi(u)=\Ee e^{iu\cdot(Y-\mu)}.
\]
Since $\|X_j\|\le2B$ and
$\sum_j\Ee\|X_j\|^2=\operatorname{tr}\Sigma=O(s)$,
\[
        \sum_j\Ee\|X_j\|^3=O(s).
\]
Consequently, uniformly for $\|u\|$ sufficiently small,
\[
        \log\Psi(u)=-\frac12u^T\Sigma u+O(s\|u\|^3).
        \tag{2.1}
\]
A symmetrization argument also gives, on a fixed neighborhood of the origin,
\[
        |\Psi(u)|\le e^{-cs\|u\|^2}.
        \tag{2.2}
\]
Indeed, if $X_j'$ is an independent copy of $X_j$, then
$1-|\Ee e^{iu\cdot X_j}|^2=
\Ee[1-\cos(u\cdot(X_j-X_j'))]$; for small $u$ the right-hand side is
bounded below by a constant multiple of $\Var(u\cdot X_j)$, and multiplication
over $j$ yields (2.2).

Fourier inversion on $[-\pi,\pi]^d$ now proves the local expansion.  Choose
$L_s\to\infty$ with $L_s^3/\sqrt{s}\to0$ and split the integral into
\[
 \|u\|\le L_s/\sqrt{s},\qquad
 L_s/\sqrt{s}<\|u\|\le\varepsilon,
 \qquad \|u\|>\varepsilon.
\]
On the first region, (2.1), followed by the change of variables
$v=\Sigma^{1/2}u$, gives the Gaussian Fourier integral with an
$o(s^{-d/2})$ error.  The estimate (2.2) makes the second region
$o(s^{-d/2})$, and tilted aperiodicity with $\lambda=0$ makes the third
region exponentially small.  The oscillatory factor
$e^{-iu\cdot(z-\mu)}$ has modulus one, so every error is uniform in $z$.
Taking absolute values in the same inversion formula proves the concentration
bound.  This is the usual fixed-dimensional Fourier proof of the lattice local
central limit theorem; compare Petrov \cite[Chapter~VII]{Petrov1975}.

It remains to prove the Gaussian local upper bound.  Let
$h=z-\mu$ and define the centered cumulant generating function
\[
        K(\lambda)=\log\Ee e^{\lambda\cdot(Y-\mu)}.
\]
For $\|\lambda\|\le\lambda_0$, its Hessian is
$\nabla^2K(\lambda)=\Var_\lambda(Y)$ and is therefore between
$c_0sI_d$ and $C_0sI_d$.  If $\|h\|=o(s)$, strong convexity shows that the
minimum of $K(\lambda)-h\cdot\lambda$ over the ball
$\|\lambda\|\le\lambda_0/2$ is attained in its interior for all large $n$.
Thus there is a unique $\lambda_h$ in that ball satisfying
\[
        \nabla K(\lambda_h)=h,
        \qquad \|\lambda_h\|\le C\|h\|/s.
\]
The associated Legendre transform
$I(h)=\lambda_h\cdot h-K(\lambda_h)$ satisfies
\[
        I(h)\ge c\|h\|^2/s.
\]
For example, the last inequality follows by testing the variational formula
$I(h)=\sup_\lambda\{\lambda\cdot h-K(\lambda)\}$ at a sufficiently small
multiple of $h/s$ and using the tilted covariance upper bound.

Under $\Pp_{\lambda_h}$, the mean of $Y$ is $z$.  Repeating the absolute-value
Fourier estimate above, now using the two tilted hypotheses, gives
\[
        \sup_w\Pp_{\lambda_h}(Y=w)\le Cs^{-d/2}.
\]
The change-of-measure identity therefore yields
\[
 \Pp(Y=z)=e^{-I(h)}\Pp_{\lambda_h}(Y=z)
 \le Cs^{-d/2}e^{-c\|h\|^2/s}.
\]
The stated displacement range is $o(s)$ because $\gamma<1/2$, completing the
proof.
\end{proof}



% END ARCHIVED BLOCK: analytic-input

% BEGIN ARCHIVED BLOCK: regular-positive-profile-theory
\section{Regular positive profiles: fixed points and short cycles}\label{sec:regular}

Let \(\sigma_n\) be the label-to-position Luce permutation generated by independent exponentials
\[
        X_{n,i}\sim \Exp(\theta_{n,i}),\qquad 1\le i\le n,
\]
so that
\[
        \sigma_n(i)=1+\#\{j:X_{n,j}<X_{n,i}\}.
\]
Throughout this section assume
\[
        \theta_{n,i}=a_n(f(i/n)+\eps_{n,i}),
        \qquad \max_{1\le i\le n}|\eps_{n,i}|\to0,
\]
where
\[
        f\in C[0,1],\qquad 0<m_f\le f(x)\le M_f<\infty.
\]
Multiplying all rates by \(a_n\) only rescales the exponential-clock time axis
and does not change their order, so take \(a_n=1\).  Fix constants
\[
        0<\underline\theta<\overline\theta<\infty
\]
such that, for every sufficiently large $n$,
$\underline\theta\le\theta_{n,i}\le\overline\theta$ for all $i$; for
example, one may take $\underline\theta=m_f/2$ and
$\overline\theta=2M_f$.

Define
\[
        A(t)=\int_0^1(1-e^{-f(u)t})\dd u,
        \qquad
        B(t)=A'(t)=\int_0^1 f(u)e^{-f(u)t}\dd u.
\]
Then
\[
        A(0)=0,\qquad A'(t)=B(t)>0,\qquad A(t)\uparrow1
        \quad(t\to\infty),
\]
so $\tau=A^{-1}:(0,1)\to(0,\infty)$ is well defined.  The density
\[
        \rho(x,y)=\frac{f(x)e^{-f(x)\tau(y)}}{B(\tau(y))},
        \qquad 0\le x\le1,\quad 0<y<1,
\]
is the Luce permuton density identified by Borga, Chatterjee, and Diaconis
\cite[Proposition~2.2]{BCD2025}.  Here it will also arise as the microscopic
point-entry density.  It is strictly positive and continuous, and is bounded
on every strip $[0,1]\times[\delta,1-\delta]$.

For \(r\ge1\) define
\[
        \lambda_r
        =\lim_{\delta\downarrow0}\frac1r
        \int_{[\delta,1-\delta]^r}
        \rho(x_1,x_2)\rho(x_2,x_3)\cdots\rho(x_r,x_1)
        \dd x_1\cdots\dd x_r,
\]
where \(x_{r+1}=x_1\).  The proof below shows that these improper limits exist and are finite.

\begin{theorem}[Regular-profile fixed points and short cycles]\label{thm:regular-cycles}
For every fixed \(R\),
\[
        (C_1(\sigma_n),\ldots,C_R(\sigma_n))
        \Rightarrow (Z_1,\ldots,Z_R),
\]
where \(Z_1,\ldots,Z_R\) are independent Poisson variables with means
\(\lambda_1,\ldots,\lambda_R\).  In particular,
\[
        C_1(\sigma_n)\Rightarrow\Pois(\lambda_1),
        \qquad
        \lambda_1=\lim_{\delta\downarrow0}\int_\delta^{1-\delta}\rho(x,x)\dd x .
\]
\end{theorem}

The proof has three parts.  First, we prove a uniform local estimate for finitely many point constraints \(\sigma_n(p_a)=q_a\) when the target positions lie in a separated bulk region.  Second, we prove endpoint and collision envelopes showing that the complement of this separated bulk region has negligible total contribution in fixed-order factorial-moment sums.  Third, we sum the separated estimate over fixed points and over disjoint directed cycles.  Thus the argument uses only fixed-order factorial moments; no global total-variation approximation is needed.

\subsection{Separated bulk local estimate}

For a finite set of specified labels \(P\subset[n]\), define empirical functions
\[
        A_{n,P}(t)=\frac1n\sum_{j\notin P}(1-e^{-\theta_{n,j}t}),
        \qquad
        B_{n,P}(t)=A_{n,P}'(t)=\frac1n\sum_{j\notin P}\theta_{n,j}e^{-\theta_{n,j}t}.
\]
If \(|P|\) is fixed, then \(A_{n,P}\to A\) and \(B_{n,P}\to B\) uniformly on compact subsets of \([0,\infty)\).

% BEGIN AUDIT REPAIR 1: uniform profile approximation
\medskip
\noindent\textit{Uniform profile approximation.}
Fix \(T<\infty\) and \(K<\infty\).  For the given triangular array satisfying
\(\max_i|\eps_{n,i}|\to0\), uniformly over all \(P\subset[n]\) with
\(|P|\le K\),
\[
 \max_{0\le m\le2}\ \sup_{0\le t\le T}
 \left|A_{n,P}^{(m)}(t)-A^{(m)}(t)\right|\longrightarrow0.
\]
Indeed, after the harmless common rescaling, the rates are uniformly bounded
above and below and converge uniformly to \(f(i/n)\); for \(m=0,1,2\) the
summands defining \(A_{n,P}^{(m)}\) are uniformly Lipschitz in the rate on
compact time intervals, removal of \(O(1)\) labels costs \(O(1/n)\), and the
remaining sums are ordinary Riemann sums.  The same argument applied to
\[
 (1-e^{-\theta t_{a\wedge b}})e^{-\theta t_{a\vee b}}
\]
and its first time derivatives shows uniform compact convergence, and hence
uniform continuity, of the normalized covariance kernels used below.  All
uniformity statements here and in the sequel are for this fixed triangular
array, uniformly over the displayed labels, bounded-size removed sets, and
admissible target positions; no quantitative rate uniform over all
perturbation arrays without a common modulus of convergence is asserted.
% END AUDIT REPAIR 1: uniform profile approximation

\begin{lemma}[Uniform saddle geometry]\label{lem:regular-saddle-geometry}
Fix \(\ell\ge1\), \(\delta>0\), and \(\eta>0\).  Uniformly over sets \(P\subset[n]\) of cardinality at most \(\ell\) and over integers
\[
        q_1<\cdots<q_\ell,
        \qquad q_a/n\in[\delta,1-\delta],
        \qquad q_{a+1}-q_a\ge \eta n,
\]
put \(y_a=(q_a-a)/n\).  For all large \(n\), the equations
\[
        A_{n,P}(t_{n,a}^*)=y_a,
        \qquad 1\le a\le\ell,
\]
have unique solutions.  These solutions lie in a compact interval
\([c_\delta,C_\delta]\subset(0,\infty)\), satisfy
\[
        t_{n,a}^*=\tau(q_a/n)+o(1),
        \qquad
        B_{n,P}(t_{n,a}^*)=B(\tau(q_a/n))+o(1),
\]
uniformly, and are separated:
\[
        t_{n,a+1}^*-t_{n,a}^*\ge c_{\delta,\eta}>0,
        \qquad 1\le a<\ell .
\]
\end{lemma}

\begin{proof}
Because \(A\) is continuous, strictly increasing, \(A(0)=0\), and \(A(t)\uparrow1\), choose \(0<c_\delta<C_\delta<\infty\) so that
\[
        A(c_\delta)<\delta/3,
        \qquad
        A(C_\delta)>1-\delta/3 .
\]
For large \(n\), uniformly over bounded-size \(P\), we have
\[
        \sup_{0\le t\le C_\delta}|A_{n,P}(t)-A(t)|\to0.
\]
Also \(y_a\in[\delta/2,1-\delta/2]\) for all large \(n\), uniformly in the allowed positions.  Hence \(A_{n,P}(c_\delta)<y_a<A_{n,P}(C_\delta)\).  Since \(|P|\le\ell\) and the rates are uniformly bounded in \([\underline\theta,\overline\theta]\), for all large \(n\),
\[
        B_{n,P}(t)
        =\frac1n\sum_{j\notin P}\theta_{n,j}e^{-\theta_{n,j}t}
        \ge \frac{\underline\theta}{2}e^{-\overline\theta t}>0
\]
on bounded intervals.  Thus \(A_{n,P}\) is strictly increasing, and the solution exists uniquely in \([c_\delta,C_\delta]\).

Uniform convergence of \(A_{n,P}\) to \(A\), together with the uniform continuity of \(\tau=A^{-1}\) on \([\delta/2,1-\delta/2]\), gives
\[
        t_{n,a}^*=\tau(y_a)+o(1)=\tau(q_a/n)+o(1),
\]
uniformly.  Uniform convergence of \(B_{n,P}\) to \(B\) on \([c_\delta,C_\delta]\) gives the displayed convergence of derivatives.

Finally, \(\tau'(y)=1/A'(\tau(y))\).  Since \(A'\) is bounded above on \([c_\delta,C_\delta]\), \(\tau'\) is bounded below on \([\delta/2,1-\delta/2]\).  Moreover,
\[
        y_{a+1}-y_a=\frac{q_{a+1}-q_a-1}{n}\ge \eta/2
\]
for all large \(n\).  Hence the limiting saddle points \(\tau(y_a)\) are separated by a positive constant depending only on \(\delta\) and \(\eta\).  The uniform convergence just proved transfers this lower bound to \(t_{n,a}^*\) for all large \(n\).
\end{proof}

\begin{lemma}[Uniform LLT for Luce cumulative counts]\label{lem:luce-cumulative-llt}
Fix \(\ell\ge1\), \(0<c<C<\infty\), \(\eta>0\), and \(K<\infty\).  Let \(P\subset[n]\) satisfy \(|P|\le K\), and let
\[
        c\le t_1<\cdots<t_\ell\le C,
        \qquad t_{a+1}-t_a\ge \eta .
\]
For independent \(X_{n,j}\sim\Exp(\theta_{n,j})\), define
\[
        S_a(t)=\sum_{j\notin P}\one_{\{X_{n,j}\le t_a\}},
        \qquad 1\le a\le\ell,
\]
and \(S(t)=(S_1(t),\ldots,S_\ell(t))\).  Then \(S(t)\) satisfies the fixed-dimensional local CLT of Lemma \ref{lem:llt} uniformly over the above parameters with scale \(s_n=n\).  In particular, if \(\Sigma_{n,P}(t)=n^{-1}\Var S(t)\), then
\[
        c_0 I_\ell\le \Sigma_{n,P}(t)\le C_0 I_\ell
\]
for constants depending only on \(\ell,c,C,\eta,K,\underline\theta,\overline\theta\), and uniformly for
\(\|z-\Ee S(t)\|=O(n^{1/2+\gamma})\), \(\gamma<1/6\),
\[
\begin{aligned}
        \Pp(S(t)=z)
        &=(2\pi n)^{-\ell/2}(\det\Sigma_{n,P}(t))^{-1/2} \\
        &\quad\times
        \exp\left[-\frac1{2n}(z-\Ee S(t))^T\Sigma_{n,P}(t)^{-1}(z-\Ee S(t))\right]
        +o(n^{-\ell/2}).
\end{aligned}
\]
Moreover, in the same range,
\[
        \Pp(S(t)=z)
        \le C n^{-\ell/2}
        \exp\left[-c\frac{\|z-\Ee S(t)\|^2}{n}\right].
\]
\end{lemma}

\begin{proof}
For a single nonspecified label $j$, the increment vector
\[
 V_j(t)=(\one_{\{X_{n,j}\le t_1\}},\ldots,
 \one_{\{X_{n,j}\le t_\ell\}})
\]
takes the $\ell+1$ adjacent cumulative values
\[
 (1,\ldots,1),\ (0,1,\ldots,1),\ldots,
 (0,\ldots,0,1),\ (0,\ldots,0).
\]
The corresponding categories are
\[
 (0,t_1],\ (t_1,t_2],\ldots,(t_{\ell-1},t_\ell],\ (t_\ell,\infty).
\]
Because the rates lie in
$[\underline\theta,\overline\theta]$, the times stay in $[c,C]$, and
successive times are separated by at least $\eta$, every category probability
is bounded below by a constant $c_*>0$, uniformly in all parameters.

The adjacent support differences are the coordinate vectors.  For any
$u\in\mathbb R^\ell$, the variance identity
\[
 \Var(u\cdot V_j)
 =\frac12\sum_{v,w}p_vp_w\{u\cdot(v-w)\}^2
\]
therefore gives
\[
        \Var(u\cdot V_j)\ge c_*^2\sum_{a=1}^\ell u_a^2.
\]
Summing over $n-|P|$ labels gives the covariance lower bound, and bounded
increments give the upper bound.

Fix $\lambda_0<\infty$.  Under any exponential tilt
$\|\lambda\|\le\lambda_0$, the probabilities of the finitely many atoms are
multiplied by factors lying between positive constants depending only on
$\ell$ and $\lambda_0$.  Thus every tilted category probability is still
bounded below uniformly, and the same covariance comparison holds under all
such tilts.  The same observation gives tilted strong aperiodicity.  Indeed,
for $u\in[-\pi,\pi]^\ell$ away from $2\pi\mathbb Z^\ell$, at least one
adjacent phase difference $u_a$ is bounded away from zero.  The two
corresponding adjacent atoms have uniformly positive tilted mass, so compactness
gives
\[
        |\Ee_\lambda e^{iu\cdot V_j(t)}|\le1-\kappa
\]
for some $\kappa>0$.  Multiplication over $n-|P|$ labels yields the exponential
Fourier bound required in Lemma \ref{lem:llt}.  All hypotheses of that lemma
are therefore satisfied with scale $s_n=n$, proving the stated local expansion
and Gaussian local upper bound.
\end{proof}

\begin{lemma}[Off-saddle exclusion]\label{lem:regular-off-saddle}
Fix \(\ell\ge1\), \(\delta>0\), \(\eta>0\), and \(\gamma\in(0,1/6)\).  In the setting of Lemma \ref{lem:regular-saddle-geometry}, let
\[
        K_\delta=[c_\delta,C_\delta]
\]
be a compact interval containing all saddles.  In the integral representation for finitely many point constraints below, the contribution from
\[
        \max_a |t_a-t_{n,a}^*|>n^{-1/2+\gamma}
\]
is \(o(n^{-\ell})\), uniformly over the allowed labels and positions.
\end{lemma}

\begin{proof}
Enlarge \(K_\delta\) slightly, if necessary, so that outside \(K_\delta\) the limiting mean \(A(t)\) is separated from \([\delta/2,1-\delta/2]\).  Uniform convergence of \(A_{n,P}\) to \(A\) then implies that if some \(t_a\notin K_\delta\), the corresponding Bernoulli sum has mean separated from its target by order \(n\).  Bernstein's inequality gives an exponentially small bound for the marginal event \(\{S_a(t)=q_a-a\}\).  The full joint point-constraint event is contained in this marginal event, so the same bound applies to the integrand.  Since the specified-clock density integrates to at most one over the full positive chamber, this part is \(O(e^{-c n})\).

It remains to consider \(t\in K_\delta^\ell\) with \(\max_a |t_a-t_{n,a}^*|>n^{-1/2+\gamma}\).  On \(K_\delta\), the derivatives \(B_{n,P}\) are uniformly bounded below.  Hence, for at least one coordinate \(a\),
\[
        |nA_{n,P}(t_a)-nA_{n,P}(t_{n,a}^*)|
        \ge c n^{1/2+\gamma}.
\]
Since \(nA_{n,P}(t_{n,a}^*)=q_a-a\), Bernstein's inequality applied to the one-dimensional count \(S_a(t)\) gives
\[
        \Pp(S_a(t)=q_a-a)\le \exp(-c n^{2\gamma}).
\]
The full joint point-constraint event is contained in this one-dimensional marginal event, so the same bound applies to the integrand.  The specified-clock density is uniformly bounded on \(K_\delta^\ell\), and the volume is bounded.  Thus the contribution of this region is exponentially small, hence \(o(n^{-\ell})\).
\end{proof}


\begin{lemma}[Gaussian-scale saddle tail]\label{lem:regular-gaussian-tail}
Fix \(\ell\ge1\), \(\delta>0\), \(\eta>0\), and \(\gamma\in(0,1/6)\).  Let \(p_1,\ldots,p_\ell\) be distinct labels and let
\(q_1<\cdots<q_\ell\) be target positions satisfying
\[
        q_a/n\in[\delta,1-\delta],\qquad q_{a+1}-q_a\ge \eta n .
\]
Let \(t_n^*=(t_{n,1}^*,\ldots,t_{n,\ell}^*)\) be the saddle vector from Lemma \ref{lem:regular-saddle-geometry}.  In the corresponding point-constraint integral, put
\[
        S(t)=\bigl(S_1(t),\ldots,S_\ell(t)\bigr),
        \qquad z_n=(q_1-1,\ldots,q_\ell-\ell).
\]
For
\[
        t=t_n^*+h/\sqrt n,\qquad \|h\|\le n^\gamma,
\]
the integrand
\[
        \prod_{a=1}^\ell\theta_{n,p_a}e^{-\theta_{n,p_a}t_a}\,
        \Pp(S(t)=z_n)
\]
is bounded by
\[
        C n^{-\ell/2}e^{-c\|h\|^2},
\]
with constants independent of the allowed labels and positions.  Consequently, uniformly in those labels and positions, the part of the saddle window with \(\|h\|>R\) contributes at most
\[
        C n^{-\ell}\int_{\|h\|>R}e^{-c\|h\|^2}\dd h .
\]
\end{lemma}

\begin{proof}
For \(\|h\|\le n^\gamma\), the saddle separation from Lemma \ref{lem:regular-saddle-geometry} implies that the coordinates of \(t\) remain ordered and separated for all large \(n\).  Moreover, \(t\) remains in a fixed compact subinterval of \((0,\infty)\).  Hence Lemma \ref{lem:luce-cumulative-llt}, including the Gaussian local upper bound inherited from Lemma \ref{lem:llt}, applies uniformly.

Let
\[
        A_n(t)=\bigl(A_{n,P}(t_1),\ldots,A_{n,P}(t_\ell)\bigr),
        \qquad z_n=(q_1-1,\ldots,q_\ell-\ell).
\]
Since \(A_n(t_n^*)=z_n/n\), Taylor's theorem and the uniform \(C^2\) bounds for \(A_{n,P}\) on the compact saddle region give
\[
        nA_n(t)-z_n
        =\sqrt n\,DA_n(t_n^*)h+O(\|h\|^2).
\]
The diagonal matrix \(DA_n(t_n^*)\) has entries \(B_{n,P}(t_{n,a}^*)\), uniformly bounded above and below.  Since \(\gamma<1/2\), the quadratic remainder is \(o(\sqrt n\|h\|)\) uniformly when \(\|h\|\le n^\gamma\) and \(\|h\|\to\infty\).  For bounded \(\|h\|\), the desired Gaussian upper bound follows from the uniform local bound \(C n^{-\ell/2}\), after decreasing \(c\).  Thus, in the range relevant for Gaussian decay,
\[
        \|nA_n(t)-z_n\|\ge c\sqrt n\|h\|.
\]
The Gaussian local upper bound gives
\[
        \Pp(S(t)=z_n)\le C n^{-\ell/2}e^{-c\|h\|^2}.
\]
The specified-clock density \(\prod_a\theta_{n,p_a}e^{-\theta_{n,p_a}t_a}\) is uniformly bounded on the compact region, which proves the integrand bound.  Since \(\dd t=n^{-\ell/2}\dd h\), the displayed tail estimate follows.
\end{proof}

\begin{lemma}[Uniform integrated lattice saddle]\label{lem:integrated-lattice-saddle}
Fix \(\ell\ge1\), \(\gamma\in(0,1/6)\), and an open set \(U\subset\Rr^\ell\).  Let \(K\Subset U\) be compact.  Suppose \(t_n^*\in K\), \(z_n\in\mathbb Z^\ell\), and \(S_n(t)\) is an integer-valued random vector such that
\[
        \Ee S_n(t)=nA_n(t),\qquad A_n(t_n^*)=z_n/n.
\]
Assume, uniformly for \(t_n^*\in K\), that:
\begin{enumerate}
\item \(A_n\) is \(C^2\) in a fixed neighborhood of \(K\), with uniformly bounded second derivatives;
\item \(D_n=DA_n(t_n^*)\) is nonsingular and \(\|D_n^{-1}\|\) is uniformly bounded;
\item the covariance matrices \(\Sigma_n(t)=n^{-1}\Var S_n(t)\) are uniformly positive definite and uniformly continuous on a fixed neighborhood of \(K\);
\item a uniform local CLT holds for \(S_n(t)\) when \(t=t_n^*+O(n^{-1/2+\gamma})\);
\item the contribution outside \(\|t-t_n^*\|\le n^{-1/2+\gamma}\) is \(o(n^{-\ell})\), and inside this window the Gaussian tail bound
\[
        |g_n(t)|\Pp(S_n(t)=z_n)\le C n^{-\ell/2}e^{-c\|h\|^2},
        \qquad t=t_n^*+h/\sqrt n,
\]
holds for some \(C,c>0\).
\end{enumerate}
If \(g_n\) is uniformly bounded and equicontinuous near \(K\), then
\[
        \int_U g_n(t)\Pp(S_n(t)=z_n)\dd t
        =n^{-\ell}g_n(t_n^*)|\det D_n|^{-1}+o(n^{-\ell}),
\]
uniformly over the family.
\end{lemma}

\begin{proof}
The off-saddle hypothesis reduces the integral to \(t=t_n^*+h/\sqrt n\) with \(\|h\|\le n^\gamma\).  By the Gaussian tail hypothesis, the part with \(\|h\|>R\) is bounded by
\[
        C n^{-\ell}\int_{\|h\|>R}e^{-c\|h\|^2}\dd h,
\]
which is \(o_R(n^{-\ell})\) uniformly as \(R\to\infty\).  It therefore suffices to evaluate the integral on each fixed ball \(\|h\|\le R\).

On such a ball, Taylor's theorem gives
\[
        z_n-nA_n(t_n^*+h/\sqrt n)
        =-\sqrt n\,D_nh+O(1),
\]
uniformly.  The covariance matrices satisfy
\[
        \Sigma_n(t_n^*+h/\sqrt n)=\Sigma_n(t_n^*)+o(1)
\]
uniformly on \(\|h\|\le R\), and \(g_n(t_n^*+h/\sqrt n)=g_n(t_n^*)+o(1)\) uniformly.  The local CLT yields
\[
\begin{aligned}
        \Pp(S_n(t_n^*+h/\sqrt n)=z_n)
        &=(2\pi n)^{-\ell/2}(\det\Sigma_n(t_n^*))^{-1/2}  \\
        &\quad\times
        \exp\left[-\frac12 h^T D_n^T\Sigma_n(t_n^*)^{-1}D_nh\right]
        +o(n^{-\ell/2}),
\end{aligned}
\]
uniformly on the fixed ball.  All \(o(1)\) terms in this fixed-ball argument are uniform over the family and over \(\|h\|\le R\).  Since \(\dd t=n^{-\ell/2}\dd h\), the contribution from \(\|h\|\le R\) equals
\[
\begin{aligned}
        &n^{-\ell}g_n(t_n^*)(2\pi)^{-\ell/2}(\det\Sigma_n(t_n^*))^{-1/2} \\
        &\qquad\times
        \int_{\|h\|\le R}
        \exp\left[-\frac12 h^T D_n^T\Sigma_n(t_n^*)^{-1}D_nh\right]\dd h
        +o_R(n^{-\ell}).
\end{aligned}
\]
Letting \(R\to\infty\), uniformly in the family, the Gaussian integral is
\[
        (2\pi)^{\ell/2}(\det\Sigma_n(t_n^*))^{1/2}|\det D_n|^{-1}.
\]
Indeed,
\[
        \det(D_n^T\Sigma_n(t_n^*)^{-1}D_n)
        =(\det D_n)^2(\det\Sigma_n(t_n^*))^{-1},
\]
so the covariance determinant cancels and the Gaussian integral contributes \(|\det D_n|^{-1}\).  Combining this with the preceding tail estimate proves the claim.
\end{proof}

\begin{proposition}[Separated bulk local estimate]\label{prop:bulk-local-final}
Fix \(\ell\ge1\), \(\delta>0\), and \(\eta>0\).  Uniformly over distinct labels \(p_1,\ldots,p_\ell\) and distinct positions \(q_1,\ldots,q_\ell\) satisfying
\[
        q_a/n\in[\delta,1-\delta],
        \qquad |q_a-q_b|\ge \eta n\quad(a\ne b),
\]
we have
\[
        \Pp(\sigma_n(p_a)=q_a,1\le a\le\ell)
        =n^{-\ell}\prod_{a=1}^\ell
        \rho(p_a/n,q_a/n)+o(n^{-\ell}).
\]
\end{proposition}

\begin{proof}
Reorder the pairs so that \(q_1<\cdots<q_\ell\), relabeling the corresponding \(p_a\)'s.  Conditional on \(X_{n,p_a}=t_a\) with \(0<t_1<\cdots<t_\ell\), define
\[
        S_a(t)=\sum_{j\notin\{p_1,\ldots,p_\ell\}}\one_{\{X_{n,j}\le t_a\}}.
\]
The event \(\sigma_n(p_a)=q_a\) for all \(a\) is equivalent to
\[
        S_a(t)=q_a-a,
        \qquad 1\le a\le\ell,
\]
because exactly \(a-1\) specified clocks occur before \(t_a\).  Thus
\[
\begin{aligned}
\Pp(\sigma_n(p)=q)
  &=\int_{0<t_1<\cdots<t_\ell}
     \prod_{a=1}^\ell\theta_{n,p_a}e^{-\theta_{n,p_a}t_a}  \\
  &\quad\times
     \Pp(S_a(t)=q_a-a,1\le a\le\ell)\dd t_1\cdots\dd t_\ell .
\end{aligned}
\]
Put \(P=\{p_1,\ldots,p_\ell\}\), \(y_a=(q_a-a)/n\), and let \(t_{n,a}^*\) be the saddle point from Lemma \ref{lem:regular-saddle-geometry}.  The ordered integration domain contains a fixed neighborhood of \(t_n^*=(t_{n,1}^*,\ldots,t_{n,\ell}^*)\), because the saddle coordinates are uniformly separated.

For fixed \(t\), the vector \(S(t)=(S_1(t),\ldots,S_\ell(t))\) has mean
\[
        \Ee S_a(t)=nA_{n,P}(t_a),
\]
and covariance \(n\Sigma_{n,P}(t)\), where
\[
        \Sigma_{n,P,ab}(t)=\frac1n\sum_{j\notin P}
        (1-e^{-\theta_{n,j}t_{a\wedge b}})e^{-\theta_{n,j}t_{a\vee b}}.
\]
In the saddle neighborhood, Lemma \ref{lem:luce-cumulative-llt} gives the required uniform local CLT.  Lemma \ref{lem:regular-gaussian-tail} supplies the Gaussian-scale tail bound, and Lemma \ref{lem:regular-off-saddle} removes the complementary region.  Apply Lemma \ref{lem:integrated-lattice-saddle} with
\[
        A_n(t)=\bigl(A_{n,P}(t_1),\ldots,A_{n,P}(t_\ell)\bigr),
        \qquad
        z_n=(q_1-1,\ldots,q_\ell-\ell),
\]
and
\[
        g_n(t)=\prod_{a=1}^\ell\theta_{n,p_a}e^{-\theta_{n,p_a}t_a}.
\]
The functions \(g_n\) are uniformly bounded and equicontinuous in the saddle neighborhood, since the specified rates are uniformly bounded and the neighborhood is compact.
% BEGIN AUDIT REPAIR 2: integrated-saddle checklist
For clarity, the hypotheses of Lemma \ref{lem:integrated-lattice-saddle}
are verified as follows.  The preceding uniform profile approximation gives
uniform \(C^2\) bounds for \(A_{n,P}\) and uniform continuity of the covariance
kernels.  Lemma \ref{lem:regular-saddle-geometry} places the saddles in a fixed
compact set and bounds the diagonal entries of \(DA_n(t_n^*)\) away from zero.
Lemma \ref{lem:luce-cumulative-llt} gives uniform covariance
nondegeneracy and the required local CLT.  Lemmas
\ref{lem:regular-off-saddle} and \ref{lem:regular-gaussian-tail} give,
respectively, the off-saddle \(o(n^{-\ell})\) contribution and the Gaussian
majorant on the saddle window.  Finally the displayed product defining
\(g_n\) is uniformly bounded and equicontinuous there.  Thus every hypothesis
of the integrated-saddle lemma is satisfied uniformly over the separated bulk
configurations.
% END AUDIT REPAIR 2: integrated-saddle checklist
Here
\[
        DA_n(t_n^*)=
        \diag(B_{n,P}(t_{n,1}^*),\ldots,B_{n,P}(t_{n,\ell}^*)).
\]
Therefore the integral equals
\[
        n^{-\ell}\prod_{a=1}^\ell
        \frac{\theta_{n,p_a}e^{-\theta_{n,p_a}t_{n,a}^*}}{B_{n,P}(t_{n,a}^*)}
        +o(n^{-\ell}).
\]
Finally, Lemma \ref{lem:regular-saddle-geometry} and the uniform convergence of the rates give
\[
        \theta_{n,p_a}=f(p_a/n)+o(1),
        \qquad
        t_{n,a}^*=\tau(q_a/n)+o(1),
        \qquad
        B_{n,P}(t_{n,a}^*)=B(\tau(q_a/n))+o(1),
\]
uniformly.  This is exactly
\[
        n^{-\ell}\prod_{a=1}^\ell
        \rho(p_a/n,q_a/n)+o(n^{-\ell}).
\]
\end{proof}

\subsection{Endpoint and collision domination}

We next prove global domination estimates sufficient for the factorial-moment sums.  Let a Luce permutation of \(N\) labels have all rates in \([\underline\theta,\overline\theta]\), and let \(R_i\) be the rank of label \(i\).  Put
\[
        \alpha=\frac{\underline\theta}{\overline\theta}\in(0,1].
\]
For \(1\le q\le N\), define
\[
        w_N(q)=
        \begin{cases}
        N^{-1}, & 1\le q\le N/2,\\[4pt]
        N^{-\alpha}(N-q+1)^{\alpha-1}, & N/2<q\le N.
        \end{cases}
\]

\begin{lemma}[One-point rank envelope]\label{lem:one-rank-envelope}
There is \(C=C(\underline\theta,\overline\theta)\) such that, for all \(N\), all rate vectors in \([\underline\theta,\overline\theta]^N\), all labels \(i\), and all \(q\),
\[
        \Pp(R_i=q)\le Cw_N(q).
\]
\end{lemma}

\begin{proof}
First suppose \(q\le N/2\).  Reveal the Luce order sequentially.  Conditional on label \(i\) not having appeared before position \(q\), at position \(q\) there are at least \(N-q+1\ge N/2\) labels remaining.  The total remaining weight is at least \(m(N-q+1)\), and the weight of label \(i\) is at most \(M\).  Hence
\[
        \Pp(R_i=q)
        \le \frac{\overline\theta}{\underline\theta(N-q+1)}
        \le \frac{2\overline\theta}{\underline\theta N}.
\]
Thus \(\Pp(R_i=q)\le C N^{-1}\) for \(q\le N/2\).

Now suppose \(q>N/2\).  Put \(s=N-q+1\).  The event \(R_i=q\) means that label \(i\) survives until there are \(s\) labels remaining and is then selected.  When there are \(r\) labels remaining and \(i\) is among them, the conditional probability that \(i\) is selected at the next draw is at least
\[
        \frac{\underline\theta}{\overline\theta r}=\frac{\alpha}{r}.
\]
Therefore the conditional probability that \(i\) is not selected at that draw is at most \(1-\alpha/r\).  Hence the probability that \(i\) survives from \(N\) labels down to \(s\) labels is at most
\[
        \prod_{r=s+1}^{N}\left(1-\frac{\alpha}{r}\right).
\]
Indeed,
\[
        \sum_{r=s+1}^{N}\frac1r\ge \log(N/s)-O(1),
\]
and therefore
\[
        \prod_{r=s+1}^{N}\left(1-\frac{\alpha}{r}\right)
        \le \exp\left(-\alpha\sum_{r=s+1}^{N}\frac1r\right)
        \le C_\alpha\left(\frac{s}{N}\right)^\alpha.
\]
Given survival to \(s\) labels, the conditional probability that \(i\) is selected next is at most
\[
        \min\left(1,\frac{\overline\theta}{\underline\theta s}\right)\le \frac{1}{\alpha s}.
\]
Therefore
\[
        \Pp(R_i=q)
        \le C\left(\frac{s}{N}\right)^\alpha\frac1s
        =C N^{-\alpha}s^{\alpha-1}.
\]
Since \(s=N-q+1\), this is the desired bound.
\end{proof}

\begin{lemma}[Gap envelope]\label{lem:gap-envelope}
Fix \(\ell\ge1\).  There is \(C_\ell\) such that the following holds.  Let \(p_1,\ldots,p_\ell\) be distinct labels and \(q_1,\ldots,q_\ell\) distinct positions in a Luce permutation of \([n]\) with rates in \([\underline\theta,\overline\theta]\).  Sort the target positions:
\[
        q_{(1)}<q_{(2)}<\cdots<q_{(\ell)}.
\]
Set \(q_{(0)}=0\), \(r_a=q_{(a)}-q_{(a-1)}\), and \(N_a=n-q_{(a-1)}\).  Then
\[
        \Pp(\sigma_n(p_a)=q_a,1\le a\le\ell)
        \le C_\ell\prod_{a=1}^\ell w_{N_a}(r_a).
\]
\end{lemma}


\begin{proof}
Let \(p_{(a)}\) denote the specified label paired with the sorted target position \(q_{(a)}\).  Reveal the Luce order from left to right, stopping at the specified target positions.  At the first specified position \(q_{(1)}\), the label \(p_{(1)}\) must appear at relative rank \(r_1=q_{(1)}\) in a Luce order of \(N_1=n\) labels.  By Lemma \ref{lem:one-rank-envelope}, this has probability at most \(Cw_{N_1}(r_1)\).

At step \(a\), condition on the full history revealed up to \(q_{(a-1)}\).  If \(p_{(a)}\) has already appeared, the conditional probability of the joint event is zero.  If \(p_{(a)}\) is still alive, the Plackett--Luce deletion property says that the unexposed tail is again a Luce order on the remaining labels, with the remaining rates.  This property follows immediately from the exponential-clock representation and memorylessness: after the revealed clocks are removed, the residual order of the unrevealed labels is generated by independent exponential clocks with their original rates.  In this tail there are \(N_a=n-q_{(a-1)}\) labels, and \(p_{(a)}\) must appear at relative rank
\[
        r_a=q_{(a)}-q_{(a-1)}.
\]
Lemma \ref{lem:one-rank-envelope} gives a conditional probability at most \(Cw_{N_a}(r_a)\).  Iterating over \(a=1,\ldots,\ell\) gives the stated bound, with \(C_\ell=C^\ell\).
\end{proof}

\begin{lemma}[Summability of the gap envelope]\label{lem:gap-sum}
There is $C<\infty$ such that
\[
        \sum_{r=1}^{N}w_N(r)\le C
\]
for every $N$.  Also, for $1\le h\le N/2$,
\[
        \sum_{r>N-h}w_N(r)\le C\left(\frac hN\right)^\alpha.
\]
For $\ell\ge1$, define
\[
H_\ell(N,h)=
\sum_{\substack{r_1,\ldots,r_\ell\ge1\\
 r_1+\cdots+r_\ell\le N\\
 N-r_1-\cdots-r_\ell\le h}}
\prod_{a=1}^\ell
w_{N-r_1-\cdots-r_{a-1}}(r_a),
\]
where the empty partial sum is zero.  Then, for every $0<\beta<\alpha$,
\[
        H_\ell(N,h)\le C_{\ell,\beta}
        \min\left\{1,\left(\frac hN\right)^\beta\right\},
        \qquad N\ge1,\ h\ge1.
\]
\end{lemma}

\begin{proof}
The first estimate follows from
\[
 \sum_{r\le N/2}N^{-1}\le\frac12
\]
and, after writing $s=N-r+1$,
\[
 \sum_{r>N/2}N^{-\alpha}(N-r+1)^{\alpha-1}
 =N^{-\alpha}\sum_{s\le N/2}s^{\alpha-1}\le C.
\]
The same substitution gives the terminal-tail estimate.

We prove the final assertion by induction on $\ell$.  If $h\ge N/2$, repeated
use of the first estimate gives $H_\ell(N,h)\le C^\ell$, which is bounded by
the asserted right-hand side after increasing the constant.  It therefore
suffices to assume $h<N/2$.

For $\ell=1$, the defining condition implies $r>N-h-1$, so the terminal-tail
bound, with a harmless endpoint term, gives
$H_1(N,h)\le C(h/N)^\alpha\le C(h/N)^\beta$.
% BEGIN AUDIT REPAIR 3A: gap endpoint term
More explicitly, the terminal-tail estimate covers \(r>N-h\).  The only
remaining possible value is \(r=N-h\); since \(h<N/2\) and
\(0<\alpha\le1\),
\[
 w_N(N-h)=N^{-\alpha}(h+1)^{\alpha-1}
 \le C N^{-\alpha}h^\alpha
 =C(h/N)^\alpha.
\]
Thus the endpoint convention does not alter the asserted power bound.
% END AUDIT REPAIR 3A: gap endpoint term
Assume the result for $\ell-1$.  After the first gap $r$, put $s=N-r$.
The induction hypothesis gives
\[
 H_\ell(N,h)
 \le C\sum_{r=1}^{N-1}w_N(r)
 \min\left\{1,\left(\frac h{N-r}\right)^\beta\right\}.
\]
For $s\le h$, the contribution is at most $C(h/N)^\alpha$.  For
$h<s<N/2$, one has $r>N/2$ and hence
\[
\begin{aligned}
 \sum_{h<s<N/2}w_N(N-s)\left(\frac hs\right)^\beta
 &\le C N^{-\alpha}h^\beta
 \sum_{h<s<N/2}s^{\alpha-1-\beta}\\
 &\le C(h/N)^\beta.
\end{aligned}
\]
Finally, for $s\ge N/2$ one has $r\le N/2$, so $w_N(r)=N^{-1}$ and each
summand is at most $CN^{-1}(h/N)^\beta$.  Summing over at most $N$ values
completes the induction.
\end{proof}

Let
\[
        E_{\delta,n}
        =\{1,\ldots,\lfloor\delta n\rfloor\}
        \cup
        \{\lceil(1-\delta)n\rceil,\ldots,n\}.
\]

\begin{lemma}[Endpoint domination]\label{lem:endpoint-dom}
Fix \(\ell\ge1\), \(0<\beta<\alpha\), and \(0<\delta<1/2\).  Then, for all sufficiently large \(n\),
\[
\sum_{\substack{q_1,
\ldots,q_\ell\ {\rm distinct}\\
q_a\in E_{\delta,n}\ {\rm for\ some}\ a}}
\sup_{p_1,
\ldots,p_\ell}
\Pp(\sigma_n(p_a)=q_a,1\le a\le\ell)
\le C_{\ell,\beta}(\delta+\delta^\beta),
\]
where \(C_{\ell,\beta}\) is independent of \(n\) and \(\delta\).
\end{lemma}

\begin{proof}
Sort the target positions.  There are at most \(\ell!\) possible orderings, so it suffices to sum over ordered target positions and use the gap envelope.

If some target position lies in the first endpoint block, then \(q_{(1)}\le\delta n\), equivalently \(r_1\le\delta n\).  For \(\delta<1/2\),
\[
        \sum_{r_1\le\delta n}w_n(r_1)
        \le \sum_{r_1\le\delta n}\frac1n\le\delta.
\]
The remaining \(\ell-1\) gaps have total envelope mass bounded by a constant by Lemma \ref{lem:gap-sum}.  Thus this endpoint block contributes at most \(C_\ell\delta\).

If some target position lies in the last endpoint block, then \(q_{(\ell)}\ge(1-\delta)n\).  Equivalently, after the last specified target position, at most \(\delta n\) labels remain.  For fixed \(\delta>0\) and all sufficiently large \(n\), we have \(1\le\delta n\le n/2\), so Lemma \ref{lem:gap-sum} applies with \(h=\delta n\).  The final estimate in that lemma gives total envelope mass at most \(C_{\ell,\beta}\delta^\beta\).  Multiplying by the finite number of orderings proves the lemma.
\end{proof}



\begin{lemma}[Bulk collision domination]\label{lem:collision-dom}
Fix \(\ell\ge1\) and \(\delta>0\).  Let \(B_{\delta,\eta,n}\) be the set of distinct \(\ell\)-tuples \(q_1,\ldots,q_\ell\) such that
\[
        q_a/n\in[\delta,1-\delta]
\]
for all \(a\), but \(|q_a-q_b|\le\eta n\) for some \(a\ne b\).  Then
\[
\sum_{q\in B_{\delta,\eta,n}}
\sup_{p_1,\ldots,p_\ell}
\Pp(\sigma_n(p_a)=q_a,1\le a\le\ell)
\le C_{\ell,\delta}\eta.
\]
\end{lemma}

\begin{proof}
Sort the positions.  If two target positions are within \(\eta n\), then some adjacent ordered pair has gap \(r_j=q_{(j)}-q_{(j-1)}\le\eta n\) with \(j\ge2\).  Since all target positions are in \([\delta n,(1-\delta)n]\), before that gap the remaining deck size satisfies \(N_j\ge\delta n\).  If \(\eta<\delta/2\), then \(r_j\le N_j/2\) and
\[
        \sum_{r_j\le\eta n}w_{N_j}(r_j)
        \le \frac{\eta n}{N_j}
        \le \frac{\eta}{\delta}.
\]
For all other gaps, we sum the gap envelope iteratively and use the uniform bound
\[
        \sum_{r=1}^N w_N(r)\le C
\]
from Lemma \ref{lem:gap-sum}.  This gives a factor depending only on \(\ell\), uniformly in the intermediate deck sizes.  Summing over the possible adjacent pair and orderings gives the result.  The case \(\eta\ge\delta/2\) is absorbed by increasing the constant.
\end{proof}

\begin{corollary}[Endpoint tightness and improper integrability]
\label{cor:endpoint-tightness}
Fix $\ell\ge1$.  For every $0<\beta<\alpha$,
\[
\lim_{\delta\downarrow0}\limsup_{n\to\infty}
\sum_{\substack{q_1,\ldots,q_\ell\ {\rm distinct}\\
q_a\in E_{\delta,n}\ {\rm for\ some}\ a}}
\sup_{p_1,\ldots,p_\ell}
\Pp(\sigma_n(p_a)=q_a,1\le a\le\ell)=0.
\]
Moreover, all limiting integrals obtained from fixed-order fixed-point and
cycle factorial moments by Proposition \ref{prop:bulk-local-final} exist as
finite improper integrals.
\end{corollary}

\begin{proof}
The displayed endpoint-tightness assertion follows immediately from Lemma
\ref{lem:endpoint-dom} by first taking $\limsup_{n\to\infty}$ and then
letting $\delta\downarrow0$.

For the integrability assertion, fix $\delta,\eta>0$.  On the truncated region
$[\delta,1-\delta]^\ell$ with pairwise target separation greater than $\eta$,
Proposition \ref{prop:bulk-local-final} identifies the relevant finite sums
with Riemann sums for the corresponding product of $\rho$-factors.  Lemma
\ref{lem:collision-dom} bounds the omitted finite-$n$ collision contribution by
$O_{\ell,\delta}(\eta)$.  Since $\rho$ is bounded and continuous on the
truncated strip, the limiting integral over the collision neighborhood also
vanishes as $\eta\downarrow0$.  Thus the full truncated integral is obtained by
first taking $n\to\infty$ and then $\eta\downarrow0$.

The gap envelope and Lemma \ref{lem:gap-sum} bound the complete fixed-order
finite sums uniformly in $n$.  Endpoint tightness shows that the portion lost
outside $[\delta,1-\delta]^\ell$ tends to zero as $\delta\downarrow0$.
Consequently the nonnegative truncated integrals are uniformly bounded and
increase to finite limits.  This applies in particular to
\[
 \int_\delta^{1-\delta}\rho(x,x)\dd x
\]
and to every cyclic trace integral
\[
 \int_{[\delta,1-\delta]^r}
 \rho(x_1,x_2)\cdots\rho(x_r,x_1)
 \dd x_1\cdots\dd x_r,
\]
as well as to their finite products in mixed factorial moments.
% BEGIN AUDIT REPAIR 3B: order of limits
The limiting procedure used here and in the factorial-moment arguments is
therefore explicit: first \(n\to\infty\) with \(\delta,\eta>0\) fixed, then
\(\eta\downarrow0\) to remove collision neighborhoods, and finally
\(\delta\downarrow0\) to remove the endpoint truncation.
% END AUDIT REPAIR 3B: order of limits
\end{proof}

\subsection{Factorial moments}

\begin{proposition}[Fixed-point factorial moments]\label{prop:regular-fixed-factorial}
For every fixed \(k\ge1\),
\[
        \Ee(C_1(\sigma_n))_k\to \lambda_1^k.
\]
\end{proposition}

\begin{proof}
Write
\[
        C_1(\sigma_n)=\sum_{i=1}^n\one_{\{\sigma_n(i)=i\}}.
\]
For fixed \(k\),
\[
        \Ee(C_1(\sigma_n))_k
        =\sum_{i_1,\ldots,i_k}^{*}
        \Pp(\sigma_n(i_a)=i_a,1\le a\le k),
\]
where \(*\) means that the indices are distinct.  Fix \(0<\delta<1/2\) and \(\eta>0\).  Split the sum into endpoint terms, non-separated bulk terms, and separated bulk terms.  Lemma \ref{lem:endpoint-dom} bounds the endpoint terms by \(O_k(\delta+\delta^\beta)\), for any fixed \(0<\beta<\alpha\).  Lemma \ref{lem:collision-dom} bounds the non-separated bulk terms by \(O_{k,\delta}(\eta)\).  On the separated bulk part, Proposition \ref{prop:bulk-local-final} gives
\[
        \Pp(\sigma_n(i_a)=i_a,1\le a\le k)
        =n^{-k}\prod_{a=1}^k\rho(i_a/n,i_a/n)+o(n^{-k}),
\]
uniformly.  The error is uniform over the separated bulk configurations, and there are \(O(n^k)\) such configurations, so its total contribution is \(o(1)\).  Hence, as \(n\to\infty\), the separated bulk contribution tends to
\[
        \left(\int_\delta^{1-\delta}\rho(x,x)\dd x\right)^k,
\]
up to an error that vanishes as \(\eta\downarrow0\).  Indeed, on \([\delta,1-\delta]\) the function \(x\mapsto\rho(x,x)\) is bounded and continuous, so the contribution of the \(\eta\)-neighborhood of the diagonals in \([\delta,1-\delta]^k\) is \(O_{k,\delta}(\eta)\).  Letting \(\eta\downarrow0\) and then \(\delta\downarrow0\), and using Corollary \ref{cor:endpoint-tightness}, gives
\[
        \Ee(C_1(\sigma_n))_k\to \lambda_1^k.
\]
\end{proof}

\begin{proposition}[Mixed short-cycle factorial moments]\label{prop:regular-cycle-factorial}
For every fixed \(R\ge1\) and fixed nonnegative integers \(k_1,\ldots,k_R\),
\[
        \Ee\prod_{r=1}^R(C_r(\sigma_n))_{k_r}
        \to \prod_{r=1}^R\lambda_r^{k_r}.
\]
\end{proposition}

\begin{proof}
Let \(\calC_{n,r}\) be the set of directed \(r\)-cycles on \([n]\), counted up to cyclic rotation.  For \(\gamma=(i_1,\ldots,i_r)\in\calC_{n,r}\), let
\[
        E_\gamma=\{\sigma_n(i_1)=i_2,
        \sigma_n(i_2)=i_3,
        \ldots,
        \sigma_n(i_r)=i_1\}.
\]
Then \(C_r(\sigma_n)=\sum_{\gamma\in\calC_{n,r}}\one_{E_\gamma}\).  Put \(L=\sum_{r=1}^R r k_r\).  Expanding the mixed factorial moment gives a sum over ordered collections of distinct cycles.  For \(k_r\) cycles of the same length \(r\), the factorial moment \((C_r(\sigma_n))_{k_r}\) keeps those cycles ordered, so no additional \(k_r!\) denominator appears; the only normalization inside a single directed \(r\)-cycle is the factor \(1/r\) for cyclic rotations.  If two distinct cycles share a label, their events are incompatible; hence only disjoint collections contribute.

% BEGIN AUDIT REPAIR 4: cycle combinatorial normalization
Equivalently, if each selected cycle is first represented by an ordered list
of its vertices, then the conversion from ordered lists to the ordered cycle
collection contributes exactly
\[
        \prod_{r=1}^R r^{-k_r}.
\]
There is no factor \(k_r!^{-1}\), because the falling factorial already orders
the \(k_r\) selected \(r\)-cycles.  Once the ordered cycle encoding is fixed,
the target tuple is determined from the source tuple by the prescribed cyclic
shifts, so only a fixed combinatorial factor depending on
\(R,k_1,\ldots,k_R\) enters the envelope estimates.
% END AUDIT REPAIR 4: cycle combinatorial normalization

For a disjoint collection, the event imposes \(L\) distinct point constraints \(\sigma_n(p_a)=q_a\).  The labels \(p_a\) are distinct, and the target positions \(q_a\) are the same labels cyclically shifted within each cycle, hence are also distinct.  The separated local estimate requires separation of the target positions; no separate metric separation of the labels is needed.

Fix \(0<\delta<1/2\) and \(\eta>0\).  Decompose the sum according to whether some target position lies in \(E_{\delta,n}\), whether all target positions lie in the bulk but two are within \(\eta n\), or whether the target positions are separated in the bulk.  For fixed \(R,k_1,\ldots,k_R\), the number of ways to arrange a given ordered \(L\)-tuple of distinct point constraints into the prescribed ordered cycle collection is bounded by a constant depending only on \(R\) and the \(k_r\)'s.  Therefore Lemmas \ref{lem:endpoint-dom} and \ref{lem:collision-dom} bound the endpoint terms by \(O_{R,k}(\delta+\delta^\beta)\) and the non-separated bulk terms by \(O_{R,k,\delta}(\eta)\).  On separated bulk configurations, Proposition \ref{prop:bulk-local-final} applies.

Consider first a single \(r\)-cycle.  On the separated bulk, the local estimate gives the Riemann sum
\[
        \frac1r
        \int_{[\delta,1-\delta]^r}
        \rho(x_1,x_2)\rho(x_2,x_3)\cdots\rho(x_r,x_1)
        \dd x_1\cdots\dd x_r,
\]
where the factor \(1/r\) appears because an ordered \(r\)-tuple represents the same directed cycle under \(r\) cyclic rotations.  For an ordered collection of disjoint cycles, the local estimate gives
\[
        n^{-L}\prod_{\gamma}\prod_{a=1}^{|\gamma|}
        \rho(x_{\gamma,a},x_{\gamma,a+1})+o(n^{-L}),
        \qquad x_{\gamma,|\gamma|+1}=x_{\gamma,1},
\]
uniformly on the separated bulk.  The disjointness restriction only removes configurations lying on equality diagonals among the \(L\) labels.  There are \(O(n^{L-1})\) such grid configurations, and on the truncated bulk the limiting product density is bounded, so their Riemann-sum weight is \(O(n^{L-1})\cdot O(n^{-L})=O(n^{-1})\).  Thus the corresponding limiting integral factors over the cycle components.  After \(n\to\infty\) and \(\eta\downarrow0\), the contribution from separated bulk collections tends to the product of the truncated cycle means.  Finally \(\delta\downarrow0\) is justified by Corollary \ref{cor:endpoint-tightness}.  We obtain
\[
        \Ee\prod_{r=1}^R(C_r(\sigma_n))_{k_r}
        \to \prod_{r=1}^R\lambda_r^{k_r}.
\]
\end{proof}

\begin{proof}[Proof of Theorem \ref{thm:regular-cycles}]
Propositions \ref{prop:regular-fixed-factorial} and
\ref{prop:regular-cycle-factorial} give convergence of all fixed mixed
factorial moments:
\[
        \Ee\prod_{r=1}^R(C_r(\sigma_n))_{k_r}
        \to
        \prod_{r=1}^R\lambda_r^{k_r}.
\]
These are exactly the mixed factorial moments of independent Poisson variables
\(Z_r\sim\Pois(\lambda_r)\).  Since ordinary mixed moments are finite linear
combinations of mixed factorial moments, all fixed mixed ordinary moments also
converge to those of \((Z_1,\ldots,Z_R)\).  The limiting vector has finite
exponential moments and is therefore moment-determinate, so the multivariate
method of moments applies; see, for example, \cite{Billingsley1995}.
% BEGIN AUDIT REPAIR 5A: moment-method reference
This citation is only to the standard convergence criterion: the model-specific
factorial-moment estimates are proved above.
% END AUDIT REPAIR 5A: moment-method reference
Thus
\[
        (C_1(\sigma_n),\ldots,C_R(\sigma_n))
        \Rightarrow (Z_1,\ldots,Z_R),
\]
with independent \(Z_r\sim\Pois(\lambda_r)\).
\end{proof}

\begin{corollary}[Total-variation convergence for fixed cycle vectors]\label{cor:regular-cycles-tv}
For every fixed $R$,
\[
 d_{\TV}\!\left(
 \mathcal L(C_1(\sigma_n),\ldots,C_R(\sigma_n)),
 \bigotimes_{r=1}^R\Pois(\lambda_r)
 \right)\longrightarrow0.
\]
\end{corollary}

\begin{proof}
Let $p_n$ and $p$ be the probability mass functions of the two laws on the
countable discrete space $\Zz_{\ge0}^R$.  Theorem \ref{thm:regular-cycles} implies
\(p_n(z)\to p(z)\) for every \(z\in\Zz_{\ge0}^R\).  This is the discrete
counting-measure form of Scheff\'e's lemma; the short argument is included
below, and standard formulations may be found in \cite{Billingsley1995}.
% BEGIN AUDIT REPAIR 5B: discrete Scheffe reference
No unbounded-moment passage is used in this total-variation step.
% END AUDIT REPAIR 5B: discrete Scheffe reference
Given \(\varepsilon>0\), choose a
finite set \(A\subset\Zz_{\ge0}^R\) with \(p(A)>1-\varepsilon\).  Then
$\sum_{z\in A}|p_n(z)-p(z)|\to0$ and $p_n(A^c)\to p(A^c)$.  Hence
\[
 \sum_{z\in\Zz_{\ge0}^R}|p_n(z)-p(z)|
 \le \sum_{z\in A}|p_n(z)-p(z)|+p_n(A^c)+p(A^c)
 \le 3\varepsilon+o(1).
\]
Letting $\varepsilon\downarrow0$ proves the claim.
\end{proof}

\subsection{Bistochasticity and long trace universality}\label{subsec:long-traces}

The fixed-length cycle means admit an operator interpretation that also determines
their large-length behavior.  Define the integral operator $K$ on
$L^2([0,1])$ by
\[
        (Kh)(x)=\int_0^1\rho(x,y)h(y)\dd y,
\]
and let $\Pi$ be the rank-one projection onto the constants,
\[
        (\Pi h)(x)=\int_0^1h(y)\dd y.
\]

\begin{lemma}[Bistochastic Hilbert--Schmidt kernel]\label{lem:bistochastic-hs}
The kernel $\rho$ is bistochastic:
\[
        \int_0^1\rho(x,y)\dd y=1
        \quad\text{for every }x\in[0,1],
        \qquad
        \int_0^1\rho(x,y)\dd x=1
        \quad\text{for every }y\in(0,1).
\]
Moreover, $K$ is a positivity-preserving Hilbert--Schmidt Markov contraction on $L^2([0,1])$ and
\[
        \|K\|_{\mathrm{HS}}^2
        =\int_{[0,1]^2}\rho(x,y)^2\dd x\dd y
        \le \frac{M_f}{m_f}.
\]
In particular, $K$ is compact, $K\one=\one$, $K^*\one=\one$, and
$K\Pi=\Pi K=\Pi$.
\end{lemma}

\begin{proof}
Put $y=A(t)$, so that $\dd y=B(t)\dd t$.  Then
\[
\begin{aligned}
        \int_0^1\rho(x,y)\dd y
        &=\int_0^\infty f(x)e^{-f(x)t}\dd t=1,\\
        \int_0^1\rho(x,y)\dd x
        &=\frac{\int_0^1f(x)e^{-f(x)\tau(y)}\dd x}
        {B(\tau(y))}=1.
\end{aligned}
\]
Thus $K\one=K^*\one=\one$.  Jensen's inequality and the two marginal
identities give
\[
        |Kh(x)|^2\le\int_0^1\rho(x,y)|h(y)|^2\dd y,
        \qquad
        \|Kh\|_2^2\le\|h\|_2^2,
\]
so $K$ is a contraction.

For the Hilbert--Schmidt norm, again use $y=A(t)$:
\[
\begin{aligned}
        \|K\|_{\mathrm{HS}}^2
        &=\int_0^\infty
        \frac{\int_0^1f(x)^2e^{-2f(x)t}\dd x}{B(t)}\dd t\\
        &\le M_f\int_0^\infty\frac{B(2t)}{B(t)}\dd t.
\end{aligned}
\]
Since $f\ge m_f$,
\[
        B(2t)=\int_0^1f(x)e^{-f(x)t}e^{-f(x)t}\dd x
        \le e^{-m_ft}B(t),
\]
and hence $\|K\|_{\mathrm{HS}}^2\le M_f/m_f$.  Hilbert--Schmidt operators are
compact.  The identities involving $\Pi$ follow from bistochasticity.
\end{proof}

\begin{lemma}[Cyclic trace formula]\label{lem:cyclic-trace-formula}
Let $K$ be a Hilbert--Schmidt integral operator on $L^2([0,1])$ with
nonnegative kernel $\rho\in L^2([0,1]^2)$.  For every $r\ge2$, $K^r$ is trace
class and
\[
 \operatorname{Tr}(K^r)
 =\int_{[0,1]^r}\rho(x_1,x_2)\rho(x_2,x_3)\cdots
 \rho(x_r,x_1)\dd x_1\cdots\dd x_r.
\]
\end{lemma}

\begin{proof}
Approximate $\rho$ in $L^2$ by nonnegative step kernels $\rho_m$, and let
$K_m$ be the corresponding finite-rank operators.  The formula is the ordinary
matrix trace identity for $K_m$.  Since $K_m\to K$ in Hilbert--Schmidt norm and
the Hilbert--Schmidt norms remain bounded, the Schatten ideal inequality
\[
        \|AB\|_1\le\|A\|_2\|B\|_2
\]
and a telescoping expansion give $K_m^r\to K^r$ in trace norm.  Hence
$\operatorname{Tr}(K_m^r)\to\operatorname{Tr}(K^r)$.  Applying the same
Cauchy--Schwarz inequalities to the cyclic products shows that their integrals
converge to the displayed integral.  This proves both its finiteness and the
identity.
\end{proof}

\begin{theorem}[Long trace universality]\label{thm:long-trace}
There are constants $C<\infty$ and $q\in(0,1)$, depending only on the limiting
profile $f$, such that, for every $r\ge2$,
\[
        \boxed{\displaystyle
        \lambda_r=\frac1r+O\!\left(\frac{q^r}{r}\right).}
\]
Equivalently,
\[
        r\lambda_r=1+O(q^r).
\]
Consequently,
\[
        \lim_{r\to\infty}\lim_{n\to\infty}
        r\,\Ee C_r(\sigma_n)=1,
        \qquad
        \sum_{r=2}^\infty\left|\lambda_r-\frac1r\right|<\infty.
\]
\end{theorem}

\begin{proof}
We first identify the peripheral spectrum of $K$.  Suppose $Kh=\zeta h$ with
$|\zeta|=1$.  Since $K$ is a contraction,
\[
        \|h\|_2=\|Kh\|_2.
\]
Equality must therefore hold in the pointwise Jensen inequality used in the
proof of Lemma \ref{lem:bistochastic-hs} for almost every $x$.  Because
$\rho(x,y)>0$ for every $x\in[0,1]$ and every $y\in(0,1)$, equality in Jensen's
inequality forces $h(y)$ to be almost everywhere constant.  It follows that
$\zeta=1$.  Thus the only eigenvectors with eigenvalue of modulus one are the
constants.  The eigenvalue $1$ is algebraically simple as well.  Indeed, if
$(K-I)g=c\one$ with $c\ne0$, then $K^ng=g+nc\one$, contradicting
$\|K^n\|\le1$.

Since $K$ is compact, its nonzero spectrum consists of isolated eigenvalues of
finite algebraic multiplicity, with possible accumulation only at zero.  Hence
there is $q\in(0,1)$ such that every eigenvalue of $K$ other than $1$ has
modulus at most $q$.  Put $T=K-\Pi$.  Then $T$ is Hilbert--Schmidt,
$T\Pi=\Pi T=0$, and
\[
        K^r=\Pi+T^r,
        \qquad r\ge1.
\]

For $r\ge2$, Lemma \ref{lem:cyclic-trace-formula} gives
\[
        \operatorname{Tr}(K^r)
        =\int_{[0,1]^r}
        \rho(x_1,x_2)\rho(x_2,x_3)\cdots\rho(x_r,x_1)
        \dd x_1\cdots\dd x_r
        =r\lambda_r.
\]
The equality with the improper integral defining $\lambda_r$ follows by
monotone convergence and Corollary \ref{cor:endpoint-tightness}.

Let $(\zeta_j)_{j\ge2}$ be the nontrivial eigenvalues of $K$, repeated according
to algebraic multiplicity.  Lidskii's theorem \cite[Theorem~3.7]{Simon2005} gives
\[
        r\lambda_r=\operatorname{Tr}(K^r)
        =1+\operatorname{Tr}(T^r)
        =1+\sum_{j\ge2}\zeta_j^r.
\]
Weyl's eigenvalue inequality for Schatten classes
\cite[Theorem~2.3]{Simon2005} gives
\[
        \sum_{j\ge2}|\zeta_j|^2\le\|T\|_{\mathrm{HS}}^2.
\]
Therefore
\[
        \left|r\lambda_r-1\right|
        \le q^{r-2}\sum_{j\ge2}|\zeta_j|^2
        \le q^{r-2}\|T\|_{\mathrm{HS}}^2.
\]
After changing the constant, this is \(O(q^r)\).
% BEGIN AUDIT REPAIR 6B: first-moment link
For each fixed \(r\), Proposition \ref{prop:regular-cycle-factorial} with
\(k_r=1\) and all other \(k_s=0\) gives
\(\Ee C_r(\sigma_n)\to\lambda_r\).  Combining this fixed-\(r\) limit with
\(r\lambda_r\to1\) just proved gives the stated iterated limit; absolute
summability follows from the same geometric bound.
% END AUDIT REPAIR 6B: first-moment link
Thus the two consequences follow.
\end{proof}


\subsection{Diagonal bands and transverse curves}\label{subsec:regular-diagonal-consequences}

The separated local estimate has several consequences for sparse deterministic
entry sets.  We first record the fixed-width diagonal-band law, and then prove a
master sparse-entry theorem.  The latter removes the transversality and
smoothness assumptions from the thin-curve statement and yields deterministic
overlap laws as a further corollary.

\begin{theorem}[Fixed-width diagonal bands]\label{thm:diagonal-band}
Under the hypotheses of Theorem \ref{thm:regular-cycles}, fix an integer
\(h\ge0\), and define
\[
        N_h=\#\{i:1\le i\le n,\ |\sigma_n(i)-i|\le h\}.
\]
Then
\[
        N_h\Rightarrow
        \Pois\left((2h+1)\lambda_1\right),
        \qquad
        \lambda_1=\lim_{\delta\downarrow0}
        \int_\delta^{1-\delta}\rho(x,x)\dd x .
\]
\end{theorem}

\begin{proof}
Write
\[
        N_h=\sum_{d=-h}^h\sum_{\substack{1\le i\le n\\1\le i+d\le n}}
        \one_{\{\sigma_n(i)=i+d\}} .
\]
Fix \(k\ge1\).  Expanding \(\Ee(N_h)_k\) gives a sum over ordered \(k\)-tuples
\((i_a,d_a)\).  Terms with repeated labels or repeated target positions are
zero unless they are identical summands, which are excluded by the factorial
expansion.  Thus the contributing terms impose \(k\) distinct point constraints
\(\sigma_n(i_a)=i_a+d_a\).

Fix \(0<\delta<1/2\) and \(\eta>0\).  Endpoint target positions contribute
\(O_{k,h}(\delta+\delta^\beta)\), for every \(0<\beta<\alpha\), by Lemma
\ref{lem:endpoint-dom}; the finite number of offsets only changes the constant.
Among the remaining bulk terms, the contribution from target positions within
\(\eta n\) of one another is \(O_{k,h,\delta}(\eta)\) by Lemma
\ref{lem:collision-dom}.  On the separated bulk part, Proposition
\ref{prop:bulk-local-final} gives, uniformly,
\[
        \Pp(\sigma_n(i_a)=i_a+d_a,\ 1\le a\le k)
        =n^{-k}\prod_{a=1}^k
        \rho(i_a/n,(i_a+d_a)/n)+o(n^{-k}).
\]
For each fixed offset \(d\), uniformly on the truncated bulk,
\[
        \rho(i/n,(i+d)/n)=\rho(i/n,i/n)+o(1),
\]
because \(\rho\) is continuous on compact subintervals of
\([0,1]\times(0,1)\).  Hence, after summing over the \((2h+1)^k\) offset
choices and passing to Riemann sums, the separated bulk contribution tends to
\[
        \left((2h+1)\int_\delta^{1-\delta}\rho(x,x)\dd x\right)^k,
\]
up to an error vanishing as \(\eta\downarrow0\).  Letting
\(\eta\downarrow0\) and then \(\delta\downarrow0\), using Corollary
\ref{cor:endpoint-tightness}, gives
\[
        \Ee(N_h)_k\to ((2h+1)\lambda_1)^k .
\]
These are the factorial moments of \(\Pois((2h+1)\lambda_1)\), which proves the
claim by the method of moments.
\end{proof}

\subsection{Sparse entry sets, thin graphs, and overlaps}\label{subsec:sparse-entry}

The preceding diagonal arguments are instances of a single sparse-entry theorem.
For a deterministic set $E_n\subset[n]^2$, put
\[
        N(E_n)=\sum_{(i,j)\in E_n}\one_{\{\sigma_n(i)=j\}},
        \qquad
        \nu_n=\frac1n\sum_{(i,j)\in E_n}
        \delta_{(i/n,j/n)}.
\]
The measures $\nu_n$ need not be probability measures.

\begin{theorem}[Sparse bulk entry sets]\label{thm:sparse-entry}
Assume the hypotheses of Theorem \ref{thm:regular-cycles}.  Fix
$\delta\in(0,1/2)$, and suppose that $E_n\subset[n]\times
[\delta n,(1-\delta)n]$ satisfies
\[
        |E_n|=O(n),
        \qquad
        \max_{1\le i\le n}\#\{j:(i,j)\in E_n\}=o(n),
\]
and
\[
        \lim_{\eta\downarrow0}\limsup_{n\to\infty}
        \frac1n\sup_{q\in[n]}
        \#\{(i,j)\in E_n:|j-q|\le\eta n\}=0.
\]
If $\nu_n\Rightarrow\nu$ weakly as finite measures on $[0,1]^2$, then
\[
        N(E_n)\Rightarrow
        \Pois\!\left(\int_{[0,1]^2}\rho(x,y)\,\nu(\dd x,\dd y)\right).
\]
\end{theorem}

\begin{proof}
Fix $k\ge1$.  Expanding the factorial moment gives
\[
        \Ee(N(E_n))_k
        =\sum_{e_1,\ldots,e_k\in E_n}^{*}
        \Pp\bigl(\sigma_n(i_a)=j_a,
        1\le a\le k\bigr),
        \qquad e_a=(i_a,j_a),
\]
where the star excludes repeated edges.  If two distinct selected edges have the
same row or the same target column, the corresponding event is impossible.

We first record a uniform bulk bound.  If $j_1,\ldots,j_k$ are distinct and all
belong to $[\delta n,(1-\delta)n]$, then Lemma
\ref{lem:gap-envelope} gives
\[
        \sup_{i_1,\ldots,i_k}
        \Pp(\sigma_n(i_a)=j_a,1\le a\le k)
        \le C_{k,\delta}n^{-k}.
\]
Indeed, after sorting the target positions, every remaining deck size is at
least $\delta n$, and every gap in Lemma \ref{lem:gap-envelope} is separated
from the terminal end of its remaining deck by at least $\delta n$; hence each
gap factor is $O_\delta(n^{-1})$.

Let
\[
        D_n=\max_i\#\{j:(i,j)\in E_n\},
        \qquad
        M_n(\eta)=\sup_q
        \#\{(i,j)\in E_n:|j-q|\le\eta n\}.
\]
The number of ordered $k$-tuples containing two distinct edges in the same row
is at most
\[
        C_k D_n|E_n|^{k-1}=o(n^k).
\]
% BEGIN AUDIT REPAIR 6C: repeated-edge bookkeeping
The unrestricted deterministic product sum used below also contains tuples
with a repeated edge.  Their number is at most
\(C_k|E_n|^{k-1}=O(n^{k-1})\), so after multiplication by the bulk
\(O(n^{-k})\) scale their contribution is \(o(1)\).  Thus excluding repeated
edges in the factorial moment and excluding repeated rows in the compatible
tuple sum are both negligible in the deterministic comparison.
% END AUDIT REPAIR 6C: repeated-edge bookkeeping
The number containing a pair of target columns at distance at most $\eta n$ is
at most
\[
        C_k |E_n|^{k-1}M_n(\eta).
\]
Consequently, by the preceding $O(n^{-k})$ probability bound, the total
factorial-moment contribution from repeated rows or $\eta n$-close target
columns is
\[
        o(1)+O_k(M_n(\eta)/n).
\]
After $n\to\infty$ and then $\eta\downarrow0$, this contribution vanishes.

On the remaining tuples the rows are distinct and the target columns are
pairwise separated by more than $\eta n$.  Proposition
\ref{prop:bulk-local-final} gives, uniformly,
\[
\begin{aligned}
        &\Pp(\sigma_n(i_a)=j_a,1\le a\le k)\\
        &\qquad=n^{-k}\prod_{a=1}^k
        \rho(i_a/n,j_a/n)+o(n^{-k}).
\end{aligned}
\]
There are $O(n^k)$ tuples, so the accumulated error is $o(1)$.  Since $\rho$ is
bounded and continuous on $[0,1]\times[\delta,1-\delta]$, the removal of the
repeated-row and close-target tuples also changes the deterministic product sum
by $o(1)$ after $n\to\infty$ and $\eta\downarrow0$.  Hence
\[
\begin{aligned}
        \Ee(N(E_n))_k
        &=\left(\frac1n\sum_{(i,j)\in E_n}
        \rho(i/n,j/n)\right)^k+o(1)\\
        &\longrightarrow
        \left(\int\rho(x,y)\,\nu(\dd x,\dd y)\right)^k.
\end{aligned}
\]
These are the factorial moments of the stated Poisson law.  The method of
moments proves the theorem.
\end{proof}

\begin{corollary}[General thin-graph Poisson law]\label{cor:thin-graph-general}
Assume the hypotheses of Theorem \ref{thm:regular-cycles}.  Let
$q_{n,i}\in[\delta n,(1-\delta)n]$ be deterministic and suppose
\[
        \frac1n\sum_{i=1}^n
        \delta_{(i/n,q_{n,i}/n)}\Rightarrow\nu,
\]
where $\nu$ is a probability measure on $[0,1]\times[\delta,1-\delta]$.
If
\[
        \lim_{\eta\downarrow0}\limsup_{n\to\infty}
        \frac1n\sup_q
        \#\{i:|q_{n,i}-q|\le\eta n\}=0,
\]
then
\[
        \#\{i:\sigma_n(i)=q_{n,i}\}
        \Rightarrow
        \Pois\!\left(\int\rho(x,y)\,\nu(\dd x,\dd y)\right).
\]
If, in addition to the preceding no-crowding assumption,
$q_{n,i}/n\to g(i/n)$ uniformly for a continuous function $g$ with values in
$[\delta,1-\delta]$, then the limiting mean is
\[
        \int_0^1\rho(x,g(x))\dd x.
\]
\end{corollary}

\begin{proof}
Apply Theorem \ref{thm:sparse-entry} to
$E_n=\{(i,q_{n,i}):1\le i\le n\}$.  Every row degree is one, and the remaining
hypotheses are exactly those in the statement.  The final assertion follows
from convergence of the empirical graph measures to the pushforward of
Lebesgue measure under $x\mapsto(x,g(x))$.
\end{proof}

\begin{corollary}[Overlaps with deterministic permutations]\label{cor:overlap-permutation}
Assume the hypotheses of Theorem \ref{thm:regular-cycles}.  Let
$\tau_n\in S_n$ be deterministic and suppose its graph permutons
\[
        \mu_{\tau_n}=\frac1n\sum_{i=1}^n
        \delta_{(i/n,\tau_n(i)/n)}
\]
converge weakly to a permuton $\mu$.  Then
\[
        \#\{i:\sigma_n(i)=\tau_n(i)\}
        \Rightarrow\Pois(\Lambda_\mu),
\]
where
\[
        \Lambda_\mu=
        \lim_{\delta\downarrow0}
        \int_{[0,1]\times[\delta,1-\delta]}
        \rho(x,y)\,\mu(\dd x,\dd y)<\infty.
\]
\end{corollary}

\begin{proof}
Fix $\delta\in(0,1/2)$ and retain only the edges
$(i,\tau_n(i))$ whose target belongs to
$[\delta n,(1-\delta)n]$.  The row degrees and column degrees are one, and
\[
        \sup_q\#\{i:|\tau_n(i)-q|\le\eta n\}
        \le 2\eta n+2.
\]
Since the second marginal of every permuton is Lebesgue measure, the boundary
lines $y=\delta$ and $y=1-\delta$ have $\mu$-measure zero.  Theorem
\ref{thm:sparse-entry} therefore gives a Poisson limit with mean
\[
        \Lambda_\mu(\delta)=
        \int_{[0,1]\times[\delta,1-\delta]}
        \rho(x,y)\,\mu(\dd x,\dd y).
\]

The expected number of omitted overlaps is bounded by
\[
        \sum_{q\in E_{\delta,n}}
        \sup_i\Pp(\sigma_n(i)=q),
\]
which tends to zero as $\delta\downarrow0$, uniformly in $n$ along the limit,
by Corollary \ref{cor:endpoint-tightness} with $\ell=1$.  Hence the omitted
count tends to zero in probability after first taking $n\to\infty$ and then
$\delta\downarrow0$.  To see that the limiting mean is finite, fix $0<\delta'<\delta$.  Applying
Theorem \ref{thm:sparse-entry} to the graph edges with target in
$[\delta',\delta]\cup[1-\delta,1-\delta']$ identifies the corresponding
limiting integral with the limit of their expected overlap count.  The latter
is bounded by the endpoint estimate above and is therefore
$O(\delta+\delta^\beta)$.  Thus \(\Lambda_\mu(\delta)\) is Cauchy as
\(\delta\downarrow0\) and has a finite monotone limit.
% BEGIN AUDIT REPAIR 5C: converging-together reference
For each fixed \(\delta\), first let \(n\to\infty\); the omitted count then
vanishes in probability as \(\delta\downarrow0\), while the truncated Poisson
means converge to \(\Lambda_\mu\).  The standard converging-together lemma
(see \cite{Billingsley1999}) therefore yields the stated Poisson limit.
% END AUDIT REPAIR 5C: converging-together reference
\end{proof}

\begin{theorem}[Transverse thin-curve Poisson law]\label{thm:thin-curve-transverse}
Assume the hypotheses of Theorem \ref{thm:regular-cycles}.  Let
$g\in C^1[0,1]$ satisfy
\[
        0<\delta_g\le g(x)\le1-\delta_g<1,
        \qquad
        \inf_{0\le x\le1}|g'(x)|>0.
\]
For all large $n$, put
\[
        q_{n,i}=\lfloor ng(i/n)\rfloor,
        \qquad
        N_g=\#\{i:\sigma_n(i)=q_{n,i}\}.
\]
Then
\[
        N_g\Rightarrow
        \Pois\left(\int_0^1\rho(x,g(x))\dd x\right).
\]
\end{theorem}

\begin{proof}
The derivative condition makes $g$ monotone and bi-Lipschitz on its image.
Hence
\[
        \sup_q\#\{i:|q_{n,i}-q|\le\eta n\}
        \le C_g(\eta n+1),
\]
while $q_{n,i}/n\to g(i/n)$ uniformly.  Corollary
\ref{cor:thin-graph-general} applies.
\end{proof}

% END ARCHIVED BLOCK: regular-positive-profile-theory

\begin{thebibliography}{99}

\bibitem{BCD2025}
J. Borga, S. Chatterjee, and P. Diaconis.
\newblock Permuton and local limits for the Luce model.
\newblock Preprint, arXiv:2509.07729v2 [math.PR], 2025.
\newblock \url{https://arxiv.org/abs/2509.07729}.
\bibitem{Billingsley1995}
P. Billingsley.
\newblock \emph{Probability and Measure}, third edition.
\newblock Wiley Series in Probability and Mathematical Statistics.
\newblock John Wiley \& Sons, New York, 1995.

\bibitem{Billingsley1999}
P. Billingsley.
\newblock \emph{Convergence of Probability Measures}, second edition.
\newblock Wiley Series in Probability and Statistics.
\newblock John Wiley \& Sons, New York, 1999.

\bibitem{Petrov1975}
V. V. Petrov.
\newblock \emph{Sums of Independent Random Variables}.
\newblock Ergebnisse der Mathematik und ihrer Grenzgebiete, Band 82.
\newblock Springer-Verlag, Berlin--Heidelberg, 1975.
\bibitem{Simon2005}
B. Simon.
\newblock \emph{Trace Ideals and Their Applications}, second edition.
\newblock Mathematical Surveys and Monographs, vol. 120.
\newblock American Mathematical Society, Providence, RI, 2005.

\end{thebibliography}

\end{document}
